\documentclass{report}
\usepackage{amsmath,amssymb}
\setlength{\parindent}{0mm}
\setlength{\parskip}{1em}
\begin{document}
\begin{center}
\rule{6in}{1pt} \
{\large Carlo Janna \\
{\bf Enhanced Block FSAI preconditioning with Domain Decomposition}}

University of Padova \\ via Trieste 63 \\ 35121 Padova \\ Italy
\\
{\tt janna@dmsa.unipd.it}\\
Massimiliano Ferronato\\
Giuseppe Gambolati\end{center}

Adaptive Block FSAI (ABF) [1,2] is a novel and effective preconditioning
technique to solve
Symmetric Positive Definite (SPD) linear systems of equations $Ax=b$ in a
parallel computing environment.
ABF is a generalization of the well-known Factored Sparse Approximate
Inverse (FSAI) [3] with
the key feature of providing a lower triangular factor $F$ such that
$FAF^T$ resembles as much as
possible a block diagonal matrix. $F$ is to some extent an "inner
preconditioner" which aims at decoupling
every diagonal block one from another. At the single processor level, ABF
may be linked to any other
scalar preconditioner, or even direct solver, which therefore performs as
an "outer preconditioner".
Recently it has been noticed that graph partitioning techniques can be
very conveniently used
to improve the ABF performance reducing both
the communication overhead and the number of iterations required for
convegence, thus greatly
enhancing the preconditioner scalability [4]. This suggests to look for a
stronger connection between ABF
and Domain Decomposition.

Consider of computing an ABF preconditioner with a set of $m$ consecutive
unknowns assigned
to the current processor $c$. The SPD matrix $A$ can be ideally partitioned as follows:
\begin{equation}
A = \left[
\begin{array}{ccc}
B & E & G \\
E^T & C & H \\
G^T & H^T & D \\
\end{array}
\right]
\end{equation}
where $\left[ \begin{array}{ccc} E^T & C & H \end{array} \right]$ includes the rows
assigned to $c$. It can be shown that the corresponding stripe of $F$ is
actually computed as:
\begin{equation}
F_i = \left[
\begin{array}{ccc}
-E^T \tilde{B}^{-1} & I_m & 0 \\
\end{array}
\right]
\end{equation}
with $\tilde{B}$ an approximation of $B$ obtained through our adaptive technique.
The $c$-th diagonal block of $FAF^T$ of size $m$, which is stored in the
current processor and is
needed for computing the outer preconditioner, takes on the following expression:
\begin{equation}
[FAF^T]_{cc} = C - E^T(2 \tilde{B}^{-1} - \tilde{B}^{-1} B \tilde{B}^{-1}
) E = C - E^T \hat{B}^{-1} E
\end{equation}
As is usually done in domain decomposition, we distinguish between
"inner" and "boundary" unknowns,
and number for each subdomain first the former and second the latter
using pedices $i_p$ and $b_p$ for
the previous processors and $i_c$ and $b_c$ for the current processor, respectively.
The matrices $E$ and $\hat{B}^{-1}$ can be written in the following block form:
\begin{equation}
E = \left[
\begin{array}{cc}
E_{i_p i_c} & E_{i_p b_c} \\
E_{b_p i_c} & E_{b_p b_c} \\
\end{array}
\right] = \left[
\begin{array}{cc}
0 & 0 \\
0 & E_{b_p b_c} \\
\end{array}
\right]
\qquad
\hat{B}^{-1} = \left[
\begin{array}{cc}
\hat{B}^{-1}_{i_p i_p} & \hat{B}^{-1}_{i_p b_p} \\
\hat{B}^{-1}_{b_p i_p} & \hat{B}^{-1}_{b_p b_p} \\
\end{array}
\right]
\label{blkEq}
\end{equation}
Three out of four blocks in $E$ are obviously zero because there is no
connection between internal
unknowns of different subdomains.
Thus the $c$-th diagonal block of $FAF^T$ becomes:
\begin{equation}
[FAF^T]_{cc} = \left[
\begin{array}{cc}
C_{i_c i_c} & C_{i_c b_c} \\
C_{i_c b_c}^T & C_{b_c b_c} - E_{b_p b_c}^T \hat{B}^{-1}_{b_p b_p} E_{b_p b_c} \\
\end{array}
\right]
\label{FAFT_cc}
\end{equation}
Equation~(\ref{FAFT_cc}) shows that ABF effects only the boundary
unknowns of the diagonal block $C$.
Then, $[FAF^T]_{cc}$ can be exactly factored and used as the outer
preconditioner for the current processor.

Assume to make a standard domain decomposition of $A$, i.e.~we exactly
factorize $C_{i_c i_c}$,
compute the Schur complement $S$ and solve an inner system with $S$.
The diagonal block of $S$ in processor $c$ reads:
\begin{equation}
S_{b_c b_c} = C_{b_c b_c} - C_{i_c b_c}^T C_{i_c i_c}^{-1} C_{i_c b_c}
\end{equation}
while the extra-diagonal blocks are equal to $E_{b_p b_c}$ and $E_{b_p
b_c}^T$ of eq.~(\ref{blkEq}).
It can be shown that solving the previous Schur complement system using
ABF as a preconditioner is equivalent
to preconditioning the whole matrix $A$ with ABF where $\tilde{B}^{-1}$
is exactly $B^{-1}$ in the computation
of $\hat{B}^{-1}_{b_p b_p}$ in~(\ref{FAFT_cc}).

Using a Schur complement domain decomposition in full leads to a more
accurate, though expensive,
approximation of $A$, while using a full matrix ABF preconditioner is
cheaper but, generally, less effective.
A good trade-off between the two approaches may be provided by using an
incomplete factorization for the "inner" blocks.
Let us number all the inner unknowns of each block first, and the
boundary unknowns second. The global
matrix $A$ can be now written as:
\begin{equation}
A = \left[
\begin{array}{cc}
A_{ii} & A_{ib} \\
A_{ib}^T & A_{bb} \\
\end{array}
\right]
\label{DDMAT}
\end{equation}
where $A_{ii}$ is a block diagonal matrix including the $C_{i_c i_c}$
blocks, $A_{ib}$ is a rectangular
matrix with the $C_{i_c b_c}$ blocks, and $A_{bb}$ is a generally sparse
matrix with the $E_{b_p b_c}$
and $C_{b_c b_c}$ blocks. $A$ in~(\ref{DDMAT}) can be preconditioned by
the incomplete factorization of
$A_{ii}$ and the computation of the approximate Schur complement $S$.
Then, the ABF preconditioner can be used to
approximately solve the inner system with $S$.

Numerical tests on large size applications show that the resulting DD-ABF
preconditioner is generally
more efficient than ABF, especially for a large number of processors,
thus ensuring a much better scalability.

\vspace{1cm}

References:

[1] C. Janna, M. Ferronato and G. Gambolati, A Block FSAI-ILU parallel preconditioner
for symmetric positive definite linear systems, SIAM Journal on
Scientific Computing, 32 (5) (2010),
pp. 2468-2484.

[2] C. Janna and M. Ferronato, Adaptive pattern research for Block FSAI preconditioners,
SIAM Journal on Scientific Computing, 33 (6) (2011), pp. 3357-3380.

[3] L. Y. Kolotilina and A. Y. Yeremin, Factorized sparse approximate
inverse preconditioning
I. Theory, SIAM J. Matrix Anal. Appl., 14 (1993), pp. 45-58.

[4] C. Janna, M. Ferronato and N. Castelletto, Graph partitioning to
improve the adaptive block
FSAI parallel preconditioner, Computer Methods in Applied Mechanics and
Engineering, submitted.


\end{document}
